\documentclass[11pt]{article}       % ana metin boyutu
\usepackage[letterpaper,                % kağıt boyutu
top=0.5in,                          % üst kenar boşluğu
bottom=0.5in,                       % alt kenar boşluğu
left=0.5in,                         % sol kenar boşluğu
right=0.5in]{geometry}              % sağ kenar boşluğu


\usepackage{comment}


% XeLaTeX için font ayarları
\usepackage{fontspec}
\setmainfont{Calibri}[Path=fonts/, UprightFont=calibri.ttf, BoldFont=calibrib.ttf,
                      ItalicFont=calibrii.ttf, BoldItalicFont=calibriz.ttf,
                      LightFont=calibril.ttf, LightItalicFont=calibrili.ttf]


%\usepackage{XCharter}               % varsayılan font yerine kullanılabilir
%\usepackage[T1]{fontenc}            % çıktı kodlaması
%\usepackage[utf8]{inputenc}         % giriş kodlaması
\usepackage{enumitem}               % madde işaretli listeler için
\usepackage[hidelinks]{hyperref}    % hyperlink biçimlendirme
\usepackage{titlesec}               % bölüm başlığı özelleştirme
\raggedright                        % metin hizalamasını kapat
\pagestyle{empty}                   % sayfa numaralandırmayı kapat


% bölüm başlıkları biçimi: kalın, boyut, üst-alt boşluk
\titleformat{\section}{\bfseries\large}{}{0pt}{}[\vspace{2pt}\titlerule\vspace{-8pt}]


% madde işaretleri biçimi
\renewcommand\labelitemi{$\vcenter{\hbox{\small$\bullet$}}$}
\setlist[itemize]{itemsep=-2pt, leftmargin=12pt, topsep=7pt}

% CV başlangıcı
\begin{document}

% isim
\centerline{\fontsize{24}{28}\selectfont \textbf{A. Tunahan Dilercan}}

\vspace{10pt}

% iletişim bilgileri
{\small
\centerline{+90 5XX XXX XX XX \,|\, \href{mailto:tunahandilercan@gmail.com}{tunahandilercan@gmail.com} \,|\, \href{https://tunahandilercan.com.tr}{tunahandilercan.com.tr}}
\vspace{3pt}
\centerline{\href{https://github.com/tunahandilercan}{github.com/tunahandilercan} \,|\, \href{https://linkedin.com/in/tunahandilercan}{linkedin.com/in/tunahandilercan}}
}


\begin{comment}

% Aşağıdaki staj deneyimleri ileride kullanılmak üzere yorum satırı olarak tutulmaktadır.

% deneyim bölümü
\section*{DENEYİM}
\textbf{UZMAR Tersanesi} \hfill \textbf{Kocaeli, TR |} Haz. 2022 -- Eyl. 2022 \\
Stajyer Mühendis \\
\vspace{-9pt}
\begin{itemize}
  \item SolidWorks'te 30+ römorkör bileşeni için parametrik şablonlar oluşturarak çizim süresini \%12 azalttı
  \item Bağlantı braketlerini el hesaplamaları ve hızlı FEA (Abaqus) ile doğrulayarak revizyon taleplerini yaklaşık \%20 düşürdü
  \item Üretim ve kalite ekipleriyle iş paketleri arasında BOM ve parça numaralandırmasını standartlaştırarak zamanında teslimat oranını yaklaşık \%15 artırdı
\end{itemize}

\textbf{Bosch} \hfill \textbf{İstanbul, TR |} Tem. 2021 -- Eyl. 2021 \\
Stajyer Mühendis \\
\vspace{-9pt}
\begin{itemize}
  \item Otomotiv bileşenleri için gömülü sistem testlerinin geliştirilmesinde görev alarak doğrulama sürecinde 12+ yazılım hatası tespit ve dokümante etti
  \item Python betikleri ile rutin veri toplama işlemlerini otomatikleştirerek manuel raporlama süresini \%40 azalttı
  \item Kıdemli mühendislerle birlikte sensör verilerini analiz ederek prototip ünitelerin performans iyileştirmelerine katkı sağladı
\end{itemize}

\textbf{Tersa Tersanesi} \hfill \textbf{Yalova, TR |} Haz. 2020 -- Ağu. 2020 \\
Stajyer Mühendis \\
\vspace{-9pt}
\begin{itemize}
  \item Mekanik tasarım ekibine SolidWorks ile gemi bileşenlerinin CAD modellerinin çizimi ve gözden geçirilmesinde destek verdi
  \item Çelik yapılar üzerinde sonlu elemanlar analizi (FEA) kontrolleri yaparak güvenlik standartlarına uygunluğu doğruladı
  \item Teknik dokümantasyon hazırlayarak ve kalite denetimlerine katkı sağlayarak zamanında raporlama doğruluğunu \%25 artırdı
\end{itemize}

\end{comment}


\vspace{-10pt}

% deneyim bölümü
\section*{DENEYİM}
\textbf{Türkiye Satranç Federasyonu} \hfill \textbf{Türkiye |} Nis. 2026 -- Devam ediyor \\
Aday Satranç Hakemi (Yarı Zamanlı) \\
\vspace{-9pt}
\begin{itemize}
  \item 2015'ten bu yana lisanslı satranç sporcusu olarak edindiği tecrübeyi profesyonel düzeye taşıyarak Türkiye Satranç Federasyonu'nun resmi eğitim ve sınavını başarıyla tamamladı
  \item Resmi turnuvalarda hakem olarak görev alarak FIDE satranç kurallarının uygulanmasını sağlıyor
\end{itemize}


\vspace{-6.5pt}

% eğitim bölümü
\section*{EĞİTİM}
\textbf{Isparta Uygulamalı Bilimler Üniversitesi} \hfill \textbf{Isparta, TR |} Beklenen Mezuniyet: 2027 \\
Mekatronik Mühendisliği (Lisans) \\[6pt]

\textbf{İstanbul Üniversitesi} \hfill \textbf{Online |} Beklenen Mezuniyet: 2028 \\
Yönetim Bilişim Sistemleri (İkinci Üniversite) \\[6pt]

\textbf{Şehit Nurullah Sabırer Fen Lisesi} \hfill \textbf{Aksaray, TR |} 2018--2022


\vspace{-6.5pt}

% projeler bölümü
\section*{PROJELER}

\textbf{Eduroid AI -- Sesli Eğitim Uygulaması} \hfill
\href{https://play.google.com/store/apps/dev?id=8414973691889516106}{\textbf{Google Play}} \\
\vspace{-9pt}
\begin{itemize}
  \item Ders ve toplantı seslerini metne dönüştüren, ardından yapay zekâ ile otomatik özet, flashcard ve test soruları üreten mobil uygulama geliştirdi
  \item Ses tanıma teknolojisi ile yapay zekâ tabanlı içerik üretimini birleştirerek öğrencilerin çalışma verimliliğini artıran uçtan uca bir çözüm sundu
\end{itemize}

\textbf{Prompt to PDF -- ChatGPT Eklentisi} \hfill
\href{https://chatgpt.com/g/g-6790c6ff33988191a78371fbc1267981-prompt-to-pdf}{\textbf{ChatGPT GPT Store}} \\
\vspace{-9pt}
\begin{itemize}
  \item LaTeX formatındaki metinleri, matematiksel formülleri ve denklemleri profesyonel kalitede PDF belgelerine dönüştüren özel bir GPT eklentisi geliştirdi
  \item Akademik ve teknik dokümanlar için doğru biçimlendirme, tutarlı tipografi ve yüksek çıktı kalitesi sağladı
\end{itemize}


\vspace{-6.5pt}

% sertifikalar bölümü
\section*{SERTİFİKALAR}
\textbf{\href{https://drive.google.com/file/d/1SBd8-eo2GcmsAc0Q561PigSK4FLun7_R/view?usp=sharing}{Siber Vatan Bootcamp}} \hfill Aralık 2023 \\
T.C. Sanayi ve Teknoloji Bakanlığı \\
\vspace{-9pt}
\begin{itemize}
  \item Sınavla seçilerek programa kabul edildi ve 176 saatlik yüz yüze teknik eğitimi başarıyla tamamladı
  \item Eğitim; Beyaz Şapkalı Hacker, Temel Linux, Sızma Testi, Zararlı Yazılım Analizi \& Tersine Mühendislik ve Uygulamalı Web Güvenliği konularını kapsadı
\end{itemize}


\vspace{-6.5pt}

% yetenekler bölümü
\section*{TEKNİK BECERİLER}
\textbf{Programlama Dilleri:} Python, Kotlin, Dart, C, Java \\
\textbf{Çatılar \& Kütüphaneler:} Flutter, Firebase, REST API, State Management \\
\textbf{Teknik İlgi Alanları:} Mobil Uygulama Geliştirme, Makine Öğrenmesi \\
\textbf{Araçlar \& Teknolojiler:} SQL, Git, Docker, Linux \\
\textbf{Yabancı Diller:} Türkçe (Ana Dil), İngilizce (B2)

\end{document}
